\subsection{Messverfahren}
\textbf{Signale 1–4}\\
Der Signalgenerator erzeugt verschiedene Signale mit unterschiedlicher Amplitude und Frequenz.
Für jedes Signal werden folgende Größen erhoben:
\begin{itemize}
    \item Periodendauer \(T\) / Frequenz \(f\)
    \item Spitze-Spitze-Spannung (\(U_\text{SP-SP}\)) und Gleichspannungsanteil \(U_\text{DC, Offset}\)
    \item Bild des Signals
\end{itemize}
Die Messungen erfolgen sowohl mit der Cursor-Funktion als auch mit den automatischen Messfunktionen, wobei der Fehler für die Cursor-Messung aus der Ablesegenauigkeit abgeschätzt wird. 

\textbf{Signal 5}\\
Signal 5 zeigt einen periodisch exponentiell abfallenden und ansteigenden Verlauf (Lade- und Entladevorgang eines Kondensators). Es wird die \textbf{Halbwertszeit} des Signals gemessen, die Zeit bis das Signal auf \(\frac{1}{2}U_\text{SP-SP}\) abfällt.
\vspace{0.5cm}\\
\textbf{qualitative Untersuchung von Signalen 6–8:}
\begin{itemize}
    \item Signal 6: Wellenpakete unterschiedlicher Höhe werden ausgesendet. Durch manuelle Triggerung wurde gezielt das größte Paket kosntant dargestellt.
    \item Signal 7: Beobachtung im Auto-Modus. Der Triggerlevel muss angepasst werden, um eine stabile Anzeige zu erhalten.
    \item Signal 8: Der Signalgenerator ist nun an zwei Kanälen des Oszis angeschlossen. Kanal 1 ist verrauscht, Kanal 2 sauber. Die Triggerung erfolgt auf Kanal 2, um beide Signale stabil darzustellen.
\end{itemize}
\vspace{0.5cm}
\textbf{Signal 9: Frequenzspektrum}\\
Signal 9 zeigt eine Schwebung, d.h. die Überlagerung zweier Sinussignale. Nach Aktivierung der AC-Kopplung und Triggerung des Signals wird die Schwingungsfrequenz \(f_1\) und die Schwebungsfrequenz \(f_2\) mittels der Cursor bestimmt. Anschließend wird die Fouriertransformation (FFT) genutzt, um die Grundfrequenzen \(f_I\) und \(f_{II}\) zu messen.
\vspace{0.5cm}\\
\textbf{Aufgabe 2 - LED}\\
Die Helligkeiten der LED wird über ein Potentiometer gedimmt, das den Tastgrad verändert. Es werden zweimal jeweils die Werte für die Pulsbreite, Pulshöhe, Periode, der Mittel- und Effektivwert bestimmt. 
\vspace{0.5cm}\\
\textbf{Aufgabe 3 - Koaxialkabel}\\
Der Signalgenerators wird über ein T-Stück mit dem Oszi verbunden,
auf die andere Seite des T-Stücks wird das Koaxialkabel angeschlossen, an dessen Ende ein kleines Gehäuse befestigt ist, das entweder als offen, kurzgeschlossen oder mit einem variablen Abschlusswiderstand eingestellt werden kann, der solange verändert wird, bis das Oszi keinen Spannungspuls mehr angezeigt.

\vfill
\subsection{Messprotokoll}
\includepdf[pages=-]{V25.pdf}